\begin{table*}[!tbp]
\centering\scriptsize\setlength\tabcolsep{2pt}\renewcommand{\arraystretch}{0.78}
\caption{Supervised concept mapping and decodability studies (decoding-family claims carrying the intrinsic-design tag are listed separately in Table~S6), part 1: probing for clinical findings and anatomy (table 1 of 5). Core medical studies only; columns and symbols as in Table~S1. Ctrl: the validity-test classes the record reports, in the codes of article Table~9 (\emph{perm} pipeline permutation control, \emph{ncal} reference-distribution calibration, \emph{ctrl} randomized control task, \emph{mobj} matched-object control, \emph{plant} planted ground truth, \emph{lrn} learned-representation attribution, \emph{int} internal-component intervention); blank means none of them is reported. The classes are not interchangeable and are never summed, so this column deliberately does not reduce to a single tick.}
\label{tab:decoding}
\begin{tabular}{@{}L{2.4cm}L{1.8cm}L{1.1cm}L{0.8cm}L{0.8cm}C{0.5cm}C{0.6cm}C{0.6cm}C{0.6cm}C{0.6cm}L{5.0cm}@{}}
\toprule
Study & Model(s) & Modality & Locus & Prov. & Ev. & Base. & Multi. & 2nd & Ctrl & Main finding \\
\midrule
\multicolumn{11}{@{}l}{\emph{Probing for clinical findings and anatomy}} \\
\cite{beyeler2024comparing} (2024) \newline {\tiny\texttt{beyeler2024comparing}} & RETFound & retina & Enc & expert & O & \xmark & \pmark & dis & -- & Regressing 17 expert retinal vascular measures on RETFound's frozen 1,024 latents (41,527 UK Biobank eyes): vascular densities R2>0.90 but tortuosity A/V ratio 0.09 and central retinal equivalents 0.06; best single latent 0.36 \\
\cite{denner2024leveraging} (2024) \newline {\tiny\texttt{denner2024leveraging}} & BiomedCLIP and other FMs & radiology & Enc & none & O & \cmark & \cmark & -- & -- & kNN retrieval over 1.6M images (P@1 up to 0.594): a specialist model beats the FMs and medical CLIP variants are not better than general ones \\
\cite{gustafsson2024evaluating} (2024) \newline {\tiny\texttt{gustafsson2024evaluating}} & 11 pathology FMs incl. UNI2-h, Virchow2 (vs Resnet-IN) & pathology & Enc & label & O & \pmark & \pmark & geo & -- & PANDA (10,616 prostate biopsy WSIs): frozen features of 11 PFMs beat an ImageNet ResNet in-distribution (UNI with ABMIL kappa 0.888+-0.013) but all collapse under Radboud->Karolinska transfer (0.247+-0.138), even for kNN \\
\cite{howard2024generative} (2024) \newline {\tiny\texttt{howard2024generative}} & CTransPath; RetCCL; UNI & pathology & Enc & label & O & \xmark & \cmark & dis & -- & GAN inversion of frozen pathology-FM embeddings (CTransPath, RetCCL, UNI; 29 cancer types) reconstructs histology preferred by pathologists in 75-94\% of comparisons; grade, subtype and gene-expression predictions from real and reconstructed tiles agree at r 0.44-0.95 \\
\cite{aerts2025foundation} (2025) \newline {\tiny\texttt{aerts2025foundation}} & ten CT FMs (FMCIB, CT-FM, CT-CLIP, Merlin, ...) & CT & Enc & label & O & \xmark & \cmark & geo & -- & kNN readouts of frozen embeddings from ten medical imaging models over six oncology cohorts (3,244 scans): FMCIB leads lung-nodule malignancy (AUC 0.886 LUNA16) while prognostic AUCs fall to 0.449-0.733 and CT-CLIP and VoCo sit near chance; no model wins everywhere; test-retest cosine 0.97-1.00 for most models (Merlin 0.81) \\
\cite{demir2025downstream} (2025) \newline {\tiny\texttt{demir2025downstream}} & SAM-Med3D; uniGradICON; SwinUNETR & MRI & Enc & label & O & \xmark & \pmark & -- & -- & Linear probes on frozen 3D bottleneck features (OAI, 2,244 patients): SAM-Med3D reaches 66.3\% KL-grade accuracy (AUC 86.6) only after affine alignment and cropping, and 84.1 AUC 24 months ahead; registration features need neither \\
\cite{germani2025bias} (2025) \newline {\tiny\texttt{germani2025bias}} & Mammo-CLIP; MedCLIP; GLoRIA (vs CLIP, DINOv2) & mammography & Enc & label & O & \cmark & \cmark & geo & -- & Linear probes on frozen features from 10 mammography datasets (26,233 scans): unified Mammo-CLIP beats GLoRIA by 15.3\%, CLIP 9.7\% and DINOv2 13.3\%; diagnosis F1 falls 0.73 internal to 0.32 external \\
\bottomrule
\end{tabular}
\end{table*}

\begin{table*}[!tbp]
\centering\scriptsize\setlength\tabcolsep{2pt}\renewcommand{\arraystretch}{0.78}
\caption{Supervised concept mapping and decodability studies (decoding-family claims carrying the intrinsic-design tag are listed separately in Table~S6), part 1: probing for clinical findings and anatomy (table 2 of 5; continued from Table~\ref{tab:decoding}). Columns and symbols as in Table~S1.}
\label{tab:decoding2}
\begin{tabular}{@{}L{2.4cm}L{1.8cm}L{1.1cm}L{0.8cm}L{0.8cm}C{0.5cm}C{0.6cm}C{0.6cm}C{0.6cm}C{0.6cm}L{5.0cm}@{}}
\toprule
Study & Model(s) & Modality & Locus & Prov. & Ev. & Base. & Multi. & 2nd & Ctrl & Main finding \\
\midrule
\multicolumn{11}{@{}l}{\emph{Probing for clinical findings and anatomy}} \\
\cite{goetze2025prediction} (2025) \newline {\tiny\texttt{goetze2025prediction}} & CT-CLIP; ClassFine; in-house Merlin-style CT VLMs & CT & Joint & label & O & \xmark & \cmark & -- & -- & Linear probes on frozen CT-VLM embeddings show lower inter-label partial correlation (0.19-0.54; PCA rank 5-11) than classification heads (0.93-0.99; rank 1): degeneracy attributed to the readout \\
\cite{kasireddy2025explainable} (2025) \newline {\tiny\texttt{kasireddy2025explainable}} & Prov-GigaPath; UNI & pathology & Enc & expert & O & \xmark & \pmark & dis & -- & 56 diabetic-nephropathy biopsies: individual Prov-GigaPath and UNI embedding dimensions correlate (|r|>=0.5) with 288 expert morphometric features, mostly colour and texture of arteries and tubules; linear probes predict two-year ESKD (AUROC 0.828 vs 0.678 for handcrafted features) \\
\cite{li2025embeddings} (2025) \newline {\tiny\texttt{li2025embeddings}} & MedImageInsight; MedSigLIP; RAD-DINO; CXR-Foundation; BiomedCLIP; Med-Flamingo & CXR & Enc & label & O & \pmark & \pmark & -- & -- & Frozen embeddings of seven encoders read by five classical adapters on 8,842 radiographs (seven tube and finding classes): MedImageInsight mAUC 93.1\%, MedSigLIP 91.0\%, RAD-DINO 90.7\%, CXR-Foundation 88.6\% vs 81.1\% for an ImageNet DenseNet-121 and 78.5\% for Med-Flamingo; spread across embeddings exceeds spread across adapters \\
\cite{li2025featurequality} (2025) \newline {\tiny\texttt{li2025featurequality}} & eight medical and general FMs & CXR & Enc & label & O & \cmark & \cmark & -- & -- & Frozen medical vs general features compared under linear probing and fine-tuning \\
\cite{marza2025thunder} (2025) \newline {\tiny\texttt{marza2025thunder}} & 23 pathology FMs & pathology & Enc & label & O & \cmark & \cmark & -- & -- & Tile-level benchmark on 16 datasets with feature-space analysis \\
\cite{neidlinger2025benchmark} (2025) \newline {\tiny\texttt{neidlinger2025benchmark}} & 19 pathology FMs & pathology & Enc & label & O & \xmark & \cmark & -- & -- & Frozen-feature benchmark on 13 external cohorts; CONCH and Virchow2 lead \\
\cite{peng2025benchmarking} (2025) \newline {\tiny\texttt{peng2025benchmarking}} & CONCH; Hibou-L; Phikon; Prov-GigaPath; UNI (vs DINOv2) & pathology & Enc & label & O & \cmark & \cmark & geo & -- & Five pathology FMs beat four general encoders on every frozen-feature placental task (UNI gestational-age MAE 2.47 weeks; region accuracy 0.829-0.895; silhouette 0.182 vs 0.079) but not under MIL, where DINOv2 competes \\
\bottomrule
\end{tabular}
\end{table*}

\begin{table*}[!tbp]
\centering\scriptsize\setlength\tabcolsep{2pt}\renewcommand{\arraystretch}{0.78}
\caption{Supervised concept mapping and decodability studies (decoding-family claims carrying the intrinsic-design tag are listed separately in Table~S6), part 1: probing for clinical findings and anatomy (table 3 of 5; continued from Table~\ref{tab:decoding}). Columns and symbols as in Table~S1.}
\label{tab:decoding3}
\begin{tabular}{@{}L{2.4cm}L{1.8cm}L{1.1cm}L{0.8cm}L{0.8cm}C{0.5cm}C{0.6cm}C{0.6cm}C{0.6cm}C{0.6cm}L{5.0cm}@{}}
\toprule
Study & Model(s) & Modality & Locus & Prov. & Ev. & Base. & Multi. & 2nd & Ctrl & Main finding \\
\midrule
\multicolumn{11}{@{}l}{\emph{Probing for clinical findings and anatomy}} \\
\cite{sadman2025interpreting} (2025) \newline {\tiny\texttt{sadman2025interpreting}} & BiomedCLIP & CXR & Joint & label & O & \xmark & \xmark & -- & -- & BiomedCLIP on IU-Xray (14 labels): linear probe on frozen embeddings macro-F1 0.183 vs zero-shot 0.105 (fine-tuning 0.235); fine-tuning raises inter/intra-class distance ratio 1.51 to 1.78 \\
\cite{tamura2025patch} (2025) \newline {\tiny\texttt{tamura2025patch}} & CONCH (vs UNI; ResNet-50) & pathology & Joint & label & O & \pmark & \cmark & dis & -- & SAE on CONCH embeddings; single neuron selected with slide labels localises tumour patches: Dice 0.61 Camelyon16 vs ABMIL 0.51; 0.57 PANDA \\
\cite{vadori2025mind} (2025) \newline {\tiny\texttt{vadori2025mind}} & UNI2; Virchow2; Prov-GigaPath (vs ImageNet-22K; LVD-142M) & pathology & Enc & label & O & \cmark & \cmark & -- & -- & Frozen multi-level patch embeddings, shared UNETR decoder, 3 datasets: general Swin2-B-22K and ConvNeXt-B-22K beat UNI2/Virchow2 by PQ +1.0 to +12.2, and shallow ViT blocks beat deep ones for every encoder \\
\cite{zhou2025generalist} (2025) \newline {\tiny\texttt{zhou2025generalist}} & retinal FMs vs general encoders & retina & Enc & label & O & \cmark & \cmark & -- & -- & Domain-specific vs generalist encoders under frozen-feature adaptation for oculomics \\
\cite{akanova2026interpreting} (2026) \newline {\tiny\texttt{akanova2026interpreting}} & EchoPrime (MViT video encoder) & echo & Enc & label & O & \xmark & \pmark & tra & -- & An MLP probe on frozen EchoPrime clip embeddings predicts LVEF at MAE 5.75\% and R2 0.631 on 88,923 Cedars-Sinai test clips (linear 7.04\%), but on 172 MIMIC-IV-Echo studies gives r=-0.094, worse than a mean baseline \\
\cite{chevli2026frozenct} (2026) \newline {\tiny\texttt{chevli2026frozenct}} & ten 3D CT encoders & CT & Enc & label & O & \xmark & \cmark & -- & -- & kNN/zero-shot/linear readouts on three cohorts; no universal winner; finding-level decodability varies \\
\cite{cui2026translating} (2026) \newline {\tiny\texttt{cui2026translating}} & CONCH; Prov-GigaPath; UNI2-h; Virchow; Virchow2 (vs ResNet50) & pathology & Enc & label & O & \pmark & \cmark & -- & -- & XGBoost on tile embeddings predicts ST-derived cell-type proportions and gene expression (leave-one-individual-out; S100A6 r=0.57 with Virchow2); combining FMs helps; ResNet50 worst \\
\bottomrule
\end{tabular}
\end{table*}

\begin{table*}[!tbp]
\centering\scriptsize\setlength\tabcolsep{2pt}\renewcommand{\arraystretch}{0.78}
\caption{Supervised concept mapping and decodability studies (decoding-family claims carrying the intrinsic-design tag are listed separately in Table~S6), part 1: probing for clinical findings and anatomy (table 4 of 5; continued from Table~\ref{tab:decoding}). Columns and symbols as in Table~S1.}
\label{tab:decoding4}
\begin{tabular}{@{}L{2.4cm}L{1.8cm}L{1.1cm}L{0.8cm}L{0.8cm}C{0.5cm}C{0.6cm}C{0.6cm}C{0.6cm}C{0.6cm}L{5.0cm}@{}}
\toprule
Study & Model(s) & Modality & Locus & Prov. & Ev. & Base. & Multi. & 2nd & Ctrl & Main finding \\
\midrule
\multicolumn{11}{@{}l}{\emph{Probing for clinical findings and anatomy}} \\
\cite{farquhar2026pretraining} (2026) \newline {\tiny\texttt{farquhar2026pretraining}} & FMCIB; CT-FM; BiomedCLIP; RAD-DINO (vs DINOv2, random init) & PET & Enc & label & O & \cmark & \cmark & geo & {\tiny lrn} & Pre-registered nine-task PET benchmark, 2,592 patients: SUV-threshold patch classification saturates at AUROC 0.97-1.00 for every frozen FM and for a 10-seed randomly initialised control (0.974); survival and subtype vary little \\
\cite{hu2026probing} (2026) \newline {\tiny\texttt{hu2026probing}} & CONCH v1.5; UNI2; Virchow; MUSK & pathology & Enc & label & O & \xmark & \cmark & -- & -- & Attention-MIL probes on four frozen pathology encoders over eight tasks in two oncology cohorts (7,634 cases) differ little across encoders but greatly across tasks (biopsy site AUC 0.92-0.94 breast vs 0.63-0.65 lung); the omics FM UCE falls below PCA; fusion helps only when no single modality carries the signal \\
\cite{jamzad2026p3ca} (2026) \newline {\tiny\texttt{jamzad2026p3ca}} & medical vision FMs & multi & Enc & none & O & \xmark & \pmark & -- & -- & Position-prompted PCA probes spatial embedding tensors encoder-agnostically \\
\cite{jung2026foundation} (2026) \newline {\tiny\texttt{jung2026foundation}} & UNI; Phikon-v2 (UNI2-h externally) & pathology & Enc & label & O & \xmark & \cmark & geo & -- & A Ridge readout on Ki-67 quantile-anchor axes of frozen UNI/Phikon-v2 slide embeddings gives R2=0.515 (95\% CI 0.380-0.633) on 117 TNBC cases and inter-encoder ICC 0.945, but external r is 0.32 and mRNA r 0.13 \\
\cite{liang2026cheap} (2026) \newline {\tiny\texttt{liang2026cheap}} & 3D CT VLMs & CT & Enc & label & O & \xmark & \pmark & -- & {\tiny perm} & Image-grounded probing benchmark predicts which training configurations succeed \\
\cite{lin2026transparent} (2026) \newline {\tiny\texttt{lin2026transparent}} & CXR contrastive FM & CXR & Enc & label & O & \xmark & \xmark & -- & -- & Linear probes against 1,074 EHR phecodes make encoder content auditable per concept \\
\cite{muthyala2026frozen} (2026) \newline {\tiny\texttt{muthyala2026frozen}} & RAD-DINO; BiomedCLIP; MedSigLIP (vs DINOv2, DINOv3) & CXR & Enc & label & O & \cmark & \cmark & -- & -- & Across five frozen ViTs and 492,724 radiographs, linear probes on pooled (CLS or patch-mean) embeddings sit at chance for lesions below about 0.4\% of image area (AUC 0.500-0.524) while patch-local pooling from the same forward pass reaches AUC >=0.99; the small-lesion gap closes only under ROI-restricted patch pooling \\
\bottomrule
\end{tabular}
\end{table*}

\begin{table*}[!tbp]
\centering\scriptsize\setlength\tabcolsep{2pt}\renewcommand{\arraystretch}{0.78}
\caption{Supervised concept mapping and decodability studies (decoding-family claims carrying the intrinsic-design tag are listed separately in Table~S6), part 1: probing for clinical findings and anatomy (table 5 of 5; continued from Table~\ref{tab:decoding}). Columns and symbols as in Table~S1.}
\label{tab:decoding5}
\begin{tabular}{@{}L{2.4cm}L{1.8cm}L{1.1cm}L{0.8cm}L{0.8cm}C{0.5cm}C{0.6cm}C{0.6cm}C{0.6cm}C{0.6cm}L{5.0cm}@{}}
\toprule
Study & Model(s) & Modality & Locus & Prov. & Ev. & Base. & Multi. & 2nd & Ctrl & Main finding \\
\midrule
\multicolumn{11}{@{}l}{\emph{Probing for clinical findings and anatomy}} \\
\cite{nooralahzadeh2026sparse} (2026) \newline {\tiny\texttt{nooralahzadeh2026sparse}} & frozen 3D CT encoders & CT & Enc & label & O & \xmark & \pmark & int & -- & Each radiological finding is carried by ~10 encoder channels; sparse channel probes match full-feature probes \\
\cite{pedersen2026look} (2026) \newline {\tiny\texttt{pedersen2026look}} & MedCLIP (ResNet-50) & CXR & Enc & label & O & \xmark & \pmark & -- & -- & 17 layer-wise probes track where CXR shortcuts become decodable across subgroups \\
\cite{ramchandani2026benchmarking} (2026) \newline {\tiny\texttt{ramchandani2026benchmarking}} & 10 pathology FMs incl. CONCH; UNI; Virchow2; PathDino & pathology & Enc & label & O & \pmark & \cmark & -- & -- & Final-block CLS attention maps of ten frozen pathology FMs classified pixel-wise by XGBoost segment tissue at Dice 0.71-0.87 (GlaS) and 0.36-0.62 (BCSS mean); all beat ImageNet ViT-B, scale does not order them \\
\cite{steiner2026when} (2026) \newline {\tiny\texttt{steiner2026when}} & TITAN; CONCH v1.5; UNIv2; H-Optimus-0; OpenMidnight & pathology & Enc & label & O & \xmark & \pmark & -- & {\tiny ncal, plant} & A label-free spurious-signal flag on frozen embeddings of five pathology FMs (8,979 TCGA patients, 16 driver genes, 50 seeds per cell) tracks post-hoc within-cancer-type AUROC collapse at Spearman 0.81-0.94, all p<0.001 \\
\cite{svatko2026deep} (2026) \newline {\tiny\texttt{svatko2026deep}} & 14 pathology FMs incl. UNIv2; CONCH; H-Optimus-1 & pathology & Enc & label & O & \cmark & \cmark & geo & {\tiny lrn} & Untrained-encoder, pixel and structure-only controls: in-domain pathology FMs lead on HEST-1k (PCC 0.37-0.68 vs 0.12-0.34 untrained, cell-centroid graphs 0.16-0.45) while untrained ViTs match them on RxRx3-core cell painting \\
\bottomrule
\end{tabular}
\end{table*}
